\section{Associations between Internal Responses, Commands, and Outcomes}
\label{sec:supp-bounded-internal}
Do larger internal responses in $\pi_{0.5}$ go with larger command changes and with lost success? Under Camera and Sensor Noise, the L18 attention and action-expert responses rank with the change in the first five commands (Table~\ref{tab:supp-internal-links}). Their links to outcomes are weak, except that the attention output under Camera separates lost from kept successes. This is an additional, post hoc study; the main paper does not rely on it.

\paragraph{Pairs.}
The analysis uses 404 of the 536 planned $\pi_{0.5}$ pairs of the internal recordings on LIBERO-Plus: the 202 settings on the 28 common tasks, each with two policy-noise draws. All 404 have the four L18 signals, verified first five commands (translation part), motion after controls 5 and 60, and a paired outcome; the other 132 pairs are not used. Pairs are matched on case, task, setting, noise draw, clean run and outcome before any signal is selected, so the sample cannot depend on the signals. Section~\ref{sec:plus-motion} reports motion on these pairs; the paired executions of Sec.~\ref{sec:five-arm} are separate episodes.

\paragraph{Signals and endpoints.}
The four signals are fixed, all at L18 and the first query: the block-input prefix state for the main camera (prefix H); the total-variation change of the attention probabilities over all keys ($P$); the attention output after projection and before gating ($W$), both recomputed; and the action-expert state from the policy's own generation (AE). $P$, $W$ and AE are taken at the first generation step. Wrist and text prefix states and the cached values V (main camera, wrist, text; no generation-step axis) are companion signals (all recording sites: Table~\ref{tab:recording-roles}). The command endpoint is the RMS of the fifteen translation differences of the first five commands whose environment call returned, in command units. The signed outcome is $dY=Y_T-Y_R$: a lost success (SF) counts $-1$, a gained success (FS) $+1$.

\paragraph{Rank correlation within tasks.}
Ranks are centered within tasks, so that differences between tasks cannot create a correlation. Within a task with $n_t$ pairs, each variable becomes $u_i=(\operatorname{rank}_{\rm avg}(x_i)-\tfrac12)/n_t$ (average ranks for ties), and pair $i$ has weight $w_i=1/(v_t n_{v(i)})$, where $v_t$ counts the task's settings and $n_{v(i)}$ the observed noise draws in pair $i$'s setting. After subtracting weighted task means, the task covariance is $C_t=\sum_iw_i\widetilde u_i\widetilde z_i$ and the variances $A_t=\sum_iw_i\widetilde u_i^2$, $B_t=\sum_iw_i\widetilde z_i^2$; the pooled coefficient is $\overline C/\sqrt{\overline A\,\overline B}$, averaging tasks equally. A variable constant within a task contributes zero variance and covariance, while the other variance still enters the denominator.

\paragraph{Lost-success AUROC and intervals.}
Among pairs whose clean run succeeds, the AUROC asks whether a larger response separates lost success (SF) from joint success (SS). Weighted SF--SS comparisons are averaged within each task, with half credit for ties, and eligible tasks are averaged equally. One shared set of 2,000 whole-task draws over the 28 tasks gives pointwise 95\% intervals; an interval is suppressed with fewer than two eligible tasks, an undefined point, or fewer than 1,900 defined draws. Signals and layer were fixed; the estimates are post hoc.

\begin{table*}[t]
\caption{\textbf{Under Camera and Sensor Noise, L18 attention and action-expert responses rank with the command change; links to outcomes are weak.} Matched $\pi_{0.5}$ LIBERO-Plus pairs: 96/102/96/110 for Background, Camera, Light and Sensor Noise (B/C/L/S). Rows: the four fixed L18 signals at the first generation step. $\rho$: task-centered weighted rank correlation of the signal with the RMS of the first five commands (translation) or with the signed outcome $dY$ (SF, lost success, $=-1$). AUROC: lost (SF) versus joint success (SS) among pairs whose clean run succeeds, over the eligible tasks in the last column; it conditions on clean success and the observed tasks. Brackets: pointwise 95\% whole-task intervals. Background and Light have only three eligible tasks; under Background, prefix H has the same score in all eligible tasks, hence the degenerate interval $[0.833,0.833]$.}

\centering\normalsize
\setlength{\tabcolsep}{3.1pt}
\begin{tabular}{@{}llcccc@{}}
\toprule
Family & L18 signal & \shortstack{$\rho$ with first\\five commands} & \shortstack{$\rho$ with\\signed $dY$} & \shortstack{AUROC:\\SF versus SS} & \shortstack{AUROC\\tasks}\\
\midrule
B & Prefix H & $-0.097[-0.397,0.186]$ & $-0.125[-0.243,0.000]$ & $0.833[0.833,0.833]$ & 3\\
B & P TV & $-0.083[-0.357,0.177]$ & $-0.081[-0.352,0.219]$ & $0.333[0.000,1.000]$ & 3\\
B & W & $0.212[-0.078,0.492]$ & $-0.178[-0.408,0.101]$ & $0.556[0.333,1.000]$ & 3\\
B & AE & $0.015[-0.282,0.296]$ & $-0.275[-0.448,-0.067]$ & $0.556[0.333,1.000]$ & 3\\
\midrule
C & Prefix H & $-0.072[-0.377,0.245]$ & $-0.185[-0.443,0.107]$ & $0.667[0.417,0.896]$ & 10\\
C & P TV & $0.348[0.075,0.596]$ & $-0.076[-0.298,0.164]$ & $0.567[0.333,0.788]$ & 10\\
C & W & $0.467[0.225,0.679]$ & $-0.284[-0.484,-0.044]$ & $0.842[0.683,0.967]$ & 10\\
C & AE & $0.400[0.116,0.657]$ & $-0.257[-0.475,0.000]$ & $0.700[0.455,0.905]$ & 10\\
\midrule
L & Prefix H & $-0.012[-0.262,0.229]$ & $-0.095[-0.233,0.000]$ & $0.611[0.500,0.833]$ & 3\\
L & P TV & $0.175[-0.117,0.468]$ & $0.143[-0.047,0.332]$ & $0.222[0.000,0.667]$ & 3\\
L & W & $-0.024[-0.344,0.294]$ & $-0.102[-0.280,0.088]$ & $0.722[0.500,1.000]$ & 3\\
L & AE & $0.000[-0.302,0.299]$ & $-0.102[-0.291,0.101]$ & $0.528[0.250,1.000]$ & 3\\
\midrule
S & Prefix H & $0.620[0.359,0.831]$ & $-0.016[-0.297,0.262]$ & $0.463[0.194,0.733]$ & 9\\
S & P TV & $0.453[0.174,0.703]$ & $-0.130[-0.360,0.099]$ & $0.593[0.333,0.867]$ & 9\\
S & W & $0.554[0.312,0.743]$ & $-0.166[-0.393,0.074]$ & $0.630[0.394,0.867]$ & 9\\
S & AE & $0.568[0.343,0.757]$ & $-0.036[-0.273,0.193]$ & $0.444[0.185,0.692]$ & 9\\
\bottomrule
\end{tabular}

\label{tab:supp-internal-links}
\end{table*}

\paragraph{What the table shows.}
All four signals rank with the command change under Sensor Noise, and $P$, $W$ and AE under Camera; under Background and Light every interval includes zero. Only two correlations with $dY$ exclude zero, Background AE and Camera $W$, both negative. Camera $W$, with AUROC $0.842[0.683,0.967]$ on 10 tasks, is the largest of the 16 AUROC estimates (Sec.~\ref{sec:auc-influence}); the other Camera and Sensor Noise intervals include one half. Camera $W$ thus has a positive command association and a negative signed-outcome association on the same pairs. It was not singled out in advance and rests on few tasks, so the main paper does not use it.

\paragraph{Attention output versus attention probabilities.}
The one contrast listed in advance, $\rho(W,\mathrm{command})-\rho(P,\mathrm{command})$, shows neither association to be stronger: all four intervals include or touch zero (B $0.295[0.000,0.592]$, C $0.119[-0.177,0.430]$, L $-0.198[-0.648,0.214]$, S $0.101[-0.124,0.336]$). The ten-condition study instead uses partial ranks, with its own endpoint and support (Sec.~\ref{supp:legacy-d3}).

\clearpage
